% 29/03/2011 at 17:00
% Written by Tordar

\documentclass{beamer}
\usepackage[brazil]{babel}		
\usepackage[utf8]{inputenc}
\usepackage{graphicx}

\renewcommand{\mathfamilydefault}{\rmdefault}

\mode<presentation>
{
  \usetheme{Warsaw} % Um de muitos temas...
  \setbeamercovered{transparent}
  \usecolortheme{dove}
  \usefonttheme{serif}
  %serif, avant, bookman, chancery, charter, euler, helvet, 
  %mathtime, mathptm, mathptmx,
	%newcent, palatino, pifont e utopia.
}
%possibilidade de usecolortheme
%albatross crane beetle dove
%fly seagull wolverine beaver

\title[Adaptação Level 2]  
{Demonstração}

\subtitle{\textcolor{blue}{UFRN}}

\author[Kumpati S. Narendra]{Kumpati S. Narendr}

\institute[Universidade Federal do Rio Grande do Norte]
{ \inst{}
  Departamento de Engenharia Elétrica\\
Universidade Federal do Rio Grande do Norte\\
}

\date[\today]{\includegraphics[width = 1 cm, height = 1.5 cm]{logoufrn.jpg}\textcolor{white}{Natal}\includegraphics[width = 1.5 cm, height = 1 cm]{ccwlogo.png} \\\small{Natal\\Junho de 2015}}

\subject{Estagio ccw}

\AtBeginSubsection[]
{
  \begin{frame}<beamer>
    \frametitle{Identificação}
    \tableofcontents[currentsection,currentsubsection]
  \end{frame}
}

\begin{document}

\begin{frame}
\titlepage
\end{frame}

\begin{frame}

\frametitle{Identificação}

\begin{enumerate}
\item<1-> Pedro Yochinori Gushiken\\
\scriptsize{Estagiário\\ Graduando em Engenharia Elétrica}

\item<2-> Local\\
\scriptsize{Empresa CCW Engenharia LTDA } 

\item<3-> Período\\
\scriptsize{Março a Junho de 2015} 

\item<4-> Wilson Medeiros de Góis\\
\scriptsize{Supervisor Técnico\\
Engenheiro Eletricista\\
Sócio Administrador da CCW Engenharia LTDA}

\item<5-> Kurios Iuri Pinheiro de Melo Queiroz\\
\scriptsize{Orientador Pedagógico\\Professor Doutor do Departamento de Engenharia Elétrica}
\end{enumerate}
\end{frame}

\begin{frame}
\frametitle{CCW Engenharia LTDA}
\begin{enumerate}
\item<1-> Fundada em 1998
\item<2-> Especializada em serviços em eletricidade em média, alta e 
em baixa tensão
\begin{enumerate}
\item<3-> \small{ ASSESSORIA E CONSULTORIA}
\item<4-> \small{PROJETO, CONSTRUÇÃO E MANUTENÇÃO }
\begin{enumerate}
\item<5-> \small{LINHAS, REDES E SUBESTAÇÕES DE ALTA E MÉDIA TENSÃO}
\item<6-> \small{MALHAS DE ATERRAMENTO}
\end{enumerate}
\end{enumerate}
\end{enumerate}
\begin{figure}
\includegraphics[width = 2.5 cm, height = 2 cm]{post.png}
\includegraphics[width = 2.5 cm, height = 2 cm]{ft_empresa.jpg}
\includegraphics[width = 2.5 cm, height = 2 cm]{seg.png}
\end{figure}
\end{frame}

\begin{frame}
\frametitle{Atividades Realizadas}
\begin{enumerate}
\item<1-> Teste Com Megômetro
\item<2-> Montagem de Quadro de Força
\item<3-> Acompanhamento da Obra do Campus de Macaíba
\end{enumerate}
\end{frame}

\begin{frame}
\frametitle{Teste com Megômetro}
\begin{columns}[t]
\begin{column}{5 cm}
Medição
\begin{enumerate}
\item<1-> Megômetro
\item<2-> Terminal Terra, Fase e Guarda
\item<3-> Tensões de Teste
\item<4-> Esquema de Ligação
\item<5-> Segurança
\end{enumerate}
\end{column}
\begin{column}{5 cm}
Procedimento
\begin{enumerate}
\item<6-> Obter Resistência de Continuidade dos Trafos
\item<7-> Repetir o teste em um intervalo razoávelmente longo
\item<8-> Observar a Evolução dos valores
\end{enumerate}
\end{column}
\end{columns}
\end{frame}


\begin{frame}
\frametitle{Esquema de Ligação}
\begin{figure}
\includegraphics[width = 10 cm, height = 7 cm]{EsquemaMegometro.png}
\caption{\scriptsize{Esquema de Ligação}}
\end{figure}
\end{frame}


\begin{frame}
\frametitle{Medidas Obtidas}
\begin{figure}
\begin{table}[htbp]
  \centering

    \begin{tabular}{cccc}
    \textbf{Trafo} & \textbf{M1} & \textbf{M2} & \textbf{M3}\\
    1     & $419.05M\Omega$ &  $412.20M\Omega$ & $408.02M\Omega$ \\
    2     & $356.46M\Omega$ &  $368.71M\Omega$   & $371.95M\Omega$  \\
    3     & $54.64M\Omega$ & $53.67M\Omega$   & $52.81M\Omega$  \\
    4     & $290.10M\Omega$& $289.11M\Omega$   & $291.64M\Omega$  \\  
    \end{tabular}%
  \caption{Medição, Março}
\end{table}

  \begin{table}[htbp]
    \centering
  
      \begin{tabular}{cccc}
      \textbf{Trafo} & \textbf{M1} & \textbf{M2} & \textbf{M3}\\
      1     & $418.02M\Omega$ & $416.1M\Omega$  & $407.89M\Omega$  \\ 
      2     & $354.14M\Omega$ & $356.66M\Omega$   & $355.51M\Omega$   \\
      3     & $69.27M\Omega$ & $60.30M\Omega$   & $57.31M\Omega$  \\  
      4     & $285.51M\Omega$ & $287.37M\Omega$   & $286.13M\Omega$  \\  
      \end{tabular}%
    \caption{Medição, Maio}
  \end{table}
  

\end{figure}
\end{frame}

\begin{frame}
\frametitle{Megômetro MG-3200}
\begin{figure}
\includegraphics[width = 6 cm, height = 7 cm]{Megometro.jpg}
\caption{\scriptsize{Foto do Megômetro utilizado}}
\end{figure}
\end{frame}


\begin{frame}
\frametitle{Montagem de Quadro de Força}
\begin{columns}[t]
\begin{column}{5 cm}
Construção
\begin{enumerate}
\item<1-> Principais materiais utilizados
\item<2-> Diagrama Unifilar
\end{enumerate}
\end{column}
\begin{column}{5 cm}
Operação
\begin{enumerate}
\item<3-> Através dos Transformadores de baixa tensão
\item<4-> Através do Gerador na Unidade de Supervisão de Corrente Alternada
\end{enumerate}
\end{column}
\end{columns}
\end{frame}


\begin{frame}
\frametitle{Principais Materiais Utilizados}
\begin{itemize}
\item 7 disjuntores BT trifásicos, $I_n=1600A, I_i = 36kA$
\item 5 disjuntores BT trifásicos, $I_n=450A, I_i = 36kA$
\item 1 disjuntor BT trifásico, $I_n = 350A, I_i = 36kA$
\item 1 disjuntor BT trifásico, $I_n = 250A, I_i = 36kA$
\item 1 chave tripolar, $I_n = 1600A$
\item 53.5 metros de barramentos chatos de cobre pintado de $2.1/2"$ por $3/8"$
\end{itemize} 
\end{frame}


\begin{frame}
\frametitle{Montagem de Quadro de Força}
\begin{figure}
\includegraphics[width = 9 cm, height = 6 cm]{base.png}
\caption{\scriptsize{Diagrama Unifilar}}
\end{figure}
\end{frame}



\begin{frame}
\frametitle{Montagem de Quadro de Força}
\begin{figure}
\includegraphics[width = 6 cm, height = 7 cm, angle=-90]{q3.jpg}
\caption{\scriptsize{Detalhe dos barramentos com seus TCs afixados}}
\end{figure}
\end{frame}

\begin{frame}
\frametitle{Montagem de Quadro de Força}
\begin{figure}
\includegraphics[width = 6 cm, height = 7 cm, angle = -90]{q2.jpg}
\caption{\scriptsize{Barramentos duplos alimentando os disjuntores}}
\end{figure}
\end{frame}


\begin{frame}
\frametitle{Operação}
\begin{figure}
\includegraphics[width = 9 cm, height = 6 cm]{normal.png}
\caption{\scriptsize{Operação Normal, através dos Transformadores de Baixa Tensão}}
\end{figure}
\end{frame}


\begin{frame}
\frametitle{Operação}
\begin{figure}
\includegraphics[width = 9 cm, height = 6 cm]{gerador.png}
\caption{\scriptsize{Alimentação das Cargas Essenciais através do Gerador na Unidade de Supervisão de Corrente Alternada}}
\end{figure}
\end{frame}



\begin{frame}
\frametitle{Acompanhamento da Obra no campus da UFRN em Macaíba}
\begin{enumerate}
\item<1-> Projeto
\begin{enumerate}
\item<2-> Rede de média tensão em cabo 336,4 AWG CAA (170$mm^2$ de secção com alma de aço)
\item<3-> Rede de baixa tensão com transformadores abaixadores distribuídos e cabo multiplexado
\item<4-> Iluminação com postes tubulares
\end{enumerate}
\item<5-> Padrão de Estruturas
\item<6-> Montagem com Munck
\end{enumerate}
\end{frame}

\begin{frame}
\frametitle{Projeto}
\begin{figure}
\includegraphics[width = 6 cm, height = 9 cm, angle=90]{ATREDE.jpg}
\caption{\scriptsize{Projeto da Rede Elétrica e da Iluminação do campus da UFRN no município de Macaíba}}
\end{figure}
\end{frame}


\begin{frame}
\frametitle{Padrão de Estruturas}
\begin{figure}
\includegraphics[width = 6 cm, height = 9 cm, angle=-90]{N1_Trafo.jpg}
\caption{\scriptsize{Estrutura N3 de final de linha com Transformador montado}}
\end{figure}
\end{frame}

\begin{frame}
\frametitle{Padrão de Estruturas}
\begin{figure}
\includegraphics[width = 6 cm, height = 9 cm, angle=-90]{N1N3.jpg}
\caption{\scriptsize{Estrutura N1-N3 (N1 tangente a linha e N3 é a cruzeta usada para derivação em $90^\circ$)}}
\end{figure}
\end{frame}


\begin{frame}
\frametitle{Padrão de Estruturas}
\begin{figure}
\includegraphics[width = 6 cm, height = 9 cm, angle=-90]{N3N3.jpg}
\caption{\scriptsize{Estrutura N3-N3 usada para "curva" fechada na linha}}
\end{figure}
\end{frame}

\begin{frame}
\frametitle{Montagem com Munck}
\begin{figure}
\includegraphics[width = 9 cm, height = 6 cm, angle=180]{Muck_Extendido.jpg}
\caption{\scriptsize{Montagem de luminária em poste tubular usando braço mecânico para erguer uma cesta com dois operários}}
\end{figure}
\end{frame}

\begin{frame}
\centering
Obrigado a Todos!
\end{frame}

\end{document}
